\section{Анализ предметной области}

\subsection{Введение в компиляторы}

Компилятор представляет собой программу, которая переводит исходный код на одном языке программирования (как правило, высокого уровня) в эквивалентный код на другом языке (чаще всего в машинный код или байт-код).

Процесс компиляции состоит из двух основных этапов: трансляция и компоновка.

Трансляция -- это основная часть компиляции, которая включает в себя анализ и преобразование исходного кода в промежуточный вид и затем в машинный код или байт-код. Трансляция программы состоит из нескольких этапов:

\begin{enumerate}
\item Лексический анализ. Последовательность символов исходного кода преобразуется в последовательность минимальных смысловых единиц - лексем.
\item Синтаксический анализ. Последовательность лексем преобразуется в синтаксическое дерево (AST), отражающее грамматическую структуру программы в соответствии с правилами языка. На этом этапе выявляются синтаксические ошибки.
\item Семантический анализ. Происходит разбор синтаксического дерева с целью установления его семантики (смысла), например, проверка области видимости переменных и функций, соответствие типов данных и т.д. Результатом данной фазы является более удобное для дальнейшей обработки промежуточное представление.
\item Оптимизация. Выполняется устранение избыточных вычислений и упрощение кода с сохранением его смысла. Оптимизация происходит на разных этапах компиляции -- например, над промежуточным кодом или над конечным машинным кодом.
\item Генерация кода. Из промежуточного представления выполняется переход в целевой код (машинный код для конкретной архитектуры процессора и операционной системы, байт-код для виртуальной машины или код на другом языке программирования). На этом этапе имеют место быть низкоуровневые оптимизации для целевой архитектуры.
\end{enumerate}

Компоновка -- процесс объединения отдельных программных модулей, полученных в процессе трансляции, в единый исполняемый файл или библиотеку. Это финальная стадия компиляции, в состав которой входит разрешение внешних ссылок между модулями, организация адресного пространств, а также оптимизации, например удаление неиспользуемого кода.

\subsection{История компиляторов}

В начале компьютерной эпохи программы писали на машинном языке или на ассемблере, учитывая архитектуру целевого процессора. Это был неудобный и трудоёмкий процесс, из-за чего появилась необходимость в инструментах для облегчения программирования.

Первым компилятором принято считать программу A-0, созданную Грейс Хоппер в начале 1950-х для компьютера UNIVAC I. A-0 позволял писать программы в виде псевдокодов, которые автоматически преобразовывались в машинные инструкции.

Также в этот период появился компилятор для языка FORTRAN, разработанный Джоном Бэкусом и его командой в IBM. FORTRAN стал одним из первых высокоуровневых языков программирования, а его компилятор -- первым широко используемым производственным компилятором.

В 1980-х разработчики компиляторов стали уделять большое внимание оптимизации целевого кода. Появились техники для увеличения скорости и уменьшения размера исполняемых программ. Также появилась потребность в портируемости программ, что привело к развитию стандартов языков и созданию промежуточных представлений (например, байт-код Java).

Современные компиляторы обладают сложной структурой и поддерживают множество языков и платформ. Развитие технологий привело к появлению компиляторов с многоуровневыми оптимизациями, динамической (Just-In-Time, JIT) компиляции -- компиляция байт-кода в машинный код во время выполнения (например JVM, CLR), проектов с открытым исходным кодом (open-source), таких как GCC и LLVM, которые предоставляют гибкие и расширяемые платформы для создания компиляторов. Кроме того, компиляторы стали интегрированы в современные среды разработки (IDE), обеспечивая не только трансляцию, но и статический анализ, автодополнение, рефакторинг и отладку кода.

\subsection{Функциональные языки программирования и Lisp}

Функциональное программирование -- парадигма программирования, которая описывает процесс выполнения программы как вычисление значений функций в математическом понимании, в отличие от императивного программирования, в котором выполнение программы определяется последовательным изменением состояний. Функциональные языки программирования характеризуются отсутствием побочных эффектов функций, выражением программ в виде композиции функций и выражений, использованием функций высшего порядка, неизменяемостью данных и ленивыми вычислениями.

Lisp -- семейство языков программирования, программы и данные в которых представляются системами односвязных списков из элементов любой природы. Это один из старейших (наряду с Fortran и Cobol) до сих пор используемых языков программирования высокого уровня, а также первый из сохранившихся в использовании языков, использующих автоматическое управление памятью и сборку мусора.

Лисп был создан Джоном Маккарти для работ по искусственному интеллекту и до сих пор остается одним из основных инструментальных средств в данной области. Применяется он и как средство обычного промышленного программирования, от встроенных скриптов до веб-приложений массового использования, хотя популярным его назвать нельзя.

Common Lisp -- диалект языка программирования Лисп, стандартизированный ANSI. Стандарт фиксирует язык как мультипарадигменный: поддерживается комбинация процедурного, функционального и объектно-ориентированного программирования. В частности, объектно-ориентированное программирование обеспечивается входящей в язык системой CLOS, а система лисп-макросов позволяет вводить в язык новые синтаксические конструкции, использовать техники метапрограммирования и обобщенного программирования.

Common Lisp отличается от таких языков, как C\#, Java, Perl, Python тем, что он определяется своим стандартом и не существует его единственной или канонической реализации. Любой желающий может ознакомиться со стандартом и создать свою собственную реализацию. Common Lisp автоматически признает все реализации как равные.

Примеры реализаций Common Lisp:

\begin{enumerate}
\item CLISP. Компилируется в байт-код, также возможна JIT-компиляция. Маленький размер образа лисп-системы. Очень эффективная длинная целочисленная арифметика. Возможность создания исполняемых файлов. FFI (интерфейс для вызова низкоуровневых функций (функций из библиотек, написанных на Си и т. п.) и для оперирования «неуправляемой» памятью). Функции обратного вызова (интеграция с «родным» кодом платформы).
\item CMUCL. Высококачественный компилятор в машинный код. FFI. Функции обратного вызова (интеграция с «родным» кодом платформы).
\item Clozure CL. Компилируется в машинный код. Быстрый компилятор. Мощный и удобный FFI. Функции обратного вызова (интеграция с бинарным кодом платформы). Возможность создания исполняемых файлов. Многопоточность на всех поддерживаемых платформах.
\item Steel Bank Common Lisp (SBCL). Компилируется в быстрый машинный код с акцентом на производительность. Мощный и стабильный FFI для вызова функций из Си и работы с памятью. Поддержка функций обратного вызова и создание исполняемых файлов. Эффективная многопоточность на большинстве платформ. Активно развивается и широко используется в промышленности.
\end{enumerate}

\subsection{Компиляция функциональных языков и Lisp}

Функциональные языки программирования традиционно ассоциируются с богатым математическим и абстрактным подходом к обработке данных. В связи с этим у них сложилась особая модель компиляции, отличающаяся от классической императивной. Функции, выражения, лямбда-выражения, которые требуют особых методов обработки при транслировании в машинный код или байт-код. Множество оптимизаций зависит от оценки выражений и замены их на более эффективные аналоги. Мета-программирование и макросистемы позволяют расширять язык, что влияет на стратегии компиляции и обработки кода.

Большинство ранних функциональных языков (например, исходные реализации Lisp в 1950–60-е) работают как интерпретаторы, выполняя код непосредственно из структур данных. Современные реализации, такие как OCaml, Haskell (GHC), Scheme (к примеру, Racket), используют компиляторы, преобразующие исходные программы в байт-код для виртуальных машин. Это обеспечивает хорошую скорость выполнения и интерактивность. Компиляция в машинный код при необходимости достигается за счёт использования специальных компиляторов (например, GHC для Haskell или SBCL для Common Lisp), которые выполняют генерацию нативного кода. В этом случае возможна высокая скорость исполнения. Ленивая (отложенная) оценка -- характерная черта многих функциональных языков, особенно Haskell, которая затрудняет традиционные методы компиляции, так как результат вычислений не всегда известен заранее.

Lisp, будучи языком с очень гибкой системой макросов и динамической типизацией, имеет свои уникальные особенности:
\begin{enumerate}
\item Интерактивное выполнение и REPL. Большинство реализаций Lisp поддерживают интерактивное программирование, где пользователь может писать код, компилировать его и запускать в реальном времени в режиме «обратной связи». Это обеспечивает быструю итерацию разработки.
\item Компиляция в байт-код и машинный код. Современные реализации Lisp позволяют компилировать программы как в байт-код (например, SBCL), так и в нативный машинный код, что значительно повышает производительность.
\item Макросы и трансформация кода. Макросы Lisp функционируют на уровне синтаксиса, позволяя писать код, который на этапе компиляции превращается в другие Lisp-выражения. Это требует специальных фаз обработки кода, что делает компиляцию более сложной, но и более мощной.
\item Динамическая типизация и гибкая модель выполнения. Это усложняет статическую оптимизацию, но современные компиляторы используют техники сборки и статического анализа для улучшения производительности.
\end{enumerate}

Современные тенденции и инструменты
\begin{enumerate}
\item Гибридные подходы. Большинство современных языков (в том числе Lisp-диалекты) используют комбинированный подход: исходный код может интерпретироваться, компилироваться в байт-код, а при необходимости -- в нативный машинный код.
\item JIT-компиляция (Just-In-Time) -- техника, при которой часть кода компилируется во время выполнения, обеспечивает баланс между быстрым запуском и высокой производительностью. Например, SBCL использует JIT для ускорения выполнения Lisp-программ.
\item Многоуровневая оптимизация. Современные компиляторы используют множество стадий оптимизации, включая статический анализ, встраивание функций, устранение избыточных вычислений и другим методам оптимизации для повышения скорости работы с функциями и представлениями данных.
\end{enumerate}

\subsection{Обзор методов оптимизации компиляторов}

Процесс оптимизации компиляторов включает в себя преобразования, которые улучшают производительность сгенерированного кода. Эти преобразования уменьшают время выполнения и используемый объем памяти без изменения смысла программы.

Появление методов оптимизации компиляторов связано с развитием вычислительных систем. Компиляторы эволюционировали от простых трансляторов к сложным инструментам с множеством оптимизаций. Исследования в этой сфере сосредоточенны преимущественно на балансе между скоростью компиляции и качеством кода. Актуальность данной темы обусловлена постоянным ростом сложности программ и аппаратного обеспечения. Современные приложения требуют эффективный код для мобильных устройств, серверов и встраиваемых систем.

\subsection{История развития методов оптимизации компиляторов}

Первые компиляторы появились в 1950-х годах. Они фокусировались на корректном переводе кода, а не на его оптимизации. В 1960-х годах разработчики осознали необходимость оптимизаций. Компилятор FORTRAN от IBM в конце 1960-х ввёл базовые техники, такие как удаление мертвого кода. Книга "The Design of an Optimizing Compiler" 1975 года описала BLISS-компилятор, который применял глобальные оптимизации.

В 1970-х годах оптимизации стали стандартом. Компиляторы для языков вроде C и Pascal включали локальные преобразования. Развитие RISC-процессоров в 1980-х годах стимулировало внедрение оптимизаций. Процессоры с конвейерной обработкой требовали код, который учитывает задержки инструкций. Компиляторы адаптировались к таким особенностям.

В 1990-х годах появились межпроцедурные оптимизации. GCC и LLVM ввели уровни оптимизации, обозначенные флагами вида -O3, активирующие наборы техник. Профильно-ориентированная оптимизация (PGO) использует данные о выполнении программы для улучшения кода. Исследования также фокусировались на параллелизме для многоядерных систем.

В 2000-х годах машинное обучение стало использоваться для оптимизации. LLVM стал платформой для экспериментов с новыми методами. Развитие мобильных устройств подчеркнуло оптимизации для энергосбережения. Техники вроде векторизации эксплуатируют SIMD-инструкции.

Современная история включает интеграцию с ИИ. Компиляторы анализируют код с помощью нейронных сетей для выбора оптимизаций. Проекты вроде TensorFlow Compiler используют специализированные методы для машинного обучения. Эволюция отражает переход от ручных настроек к автоматизированным системам.

\subsection{Основные концепции оптимизации компиляторов}

Оптимизация компиляторов основана на преобразованиях, которые сохраняют семантику программы. Преобразования применяются к промежуточному представлению (IR), такому как трёхадресный код или SSA-форма. SSA (Static Single Assignment) присваивает каждой переменной значение только один раз, что упрощает анализ.

Ключевой концепцией служит анализ потока данных. Он определяет, как значения распространяются через программу. Анализ включает живые переменные, достижимые определения и константные значения. Эти данные позволяют устранять избыточные вычисления.

Другая концепция - анализ контрольного потока. Граф контрольного потока (CFG) моделирует ветвления и циклы. Оптимизации используют CFG для выявления недостижимого кода или слияния блоков.

Профильно-ориентированная оптимизация собирает статистику выполнения. Инструменты вроде perf или gprof предоставляют данные о горячих точках. Компилятор перестраивает код на основе этих данных, перемещая редкие ветви в конец.

Концепция уровней оптимизации определяет наборы техник. Уровень -O0 отключает оптимизации для отладки. Уровень -O3 активирует агрессивные преобразования, включая векторизацию и встраивание функций.

Оптимизации учитывают компромиссы. Улучшение скорости может увеличить размер кода. Оптимизации для размера (-Os) минимизируют бинарный файл. Концепция безопасности гарантирует, что преобразования не вводят ошибки.

В функциональных языках оптимизации включают хвостовую рекурсию. Это преобразует рекурсию в цикл для экономии стека.

\subsection{Классификация методов оптимизации}

Методы оптимизации классифицируют по нескольким критериям:

\begin{enumerate}
\item В зависимости от машины. Независимые от машины методы работают с промежуточным кодом без учета архитектуры. Они включают константное сворачивание и удаление мертвого кода. Зависимые от машины методы адаптируют код к процессору, включая выбор инструкций и планирование.
\item По области применения. Локальные оптимизации ограничиваются базовым блоком - последовательностью инструкций без ветвлений. Они требуют минимального анализа. Глобальные оптимизации охватывают функцию, анализируя поток данных через блоки. Межпроцедурные оптимизации рассматривают всю программу, включая инлайнинг функций.
\item По типу преобразования. Оптимизации циклов фокусируются на петлях, которые занимают основное время выполнения. Техники включают развертку циклов и перемещение инвариантов. Оптимизации на основе потока данных устраняют избыточные выражения. SSA-оптимизации используют форму SSA для глобальной нумерации значений.
\item По фазе компиляции. Оптимизации фронтенда обрабатывают исходный код, включая макросы. Бэкенд-оптимизации генерируют машинный код, включая аллокацию регистров.
\end{enumerate}

\subsection{Обзор основных методов оптимизации}

\begin{enumerate}
\item Константное сворачивание. Вычисляет выражения с константами на этапе компиляции. Пример: выражение 2 + 3 заменяется на 5. Это сокращает инструкции и уменьшает время выполнения.
\item Пропагация констант. Заменяет переменные их значениями при вычислении выражений, в которых используются эти переменные. Если переменная a равна 10, то a * 2 становится 20. Техника сочетается со сворачиванием.
\item Удаление мертвого кода. Устраняет инструкции без влияния на результат. Анализ живых переменных выявляет такие инструкции. Пример: присваивание переменной, которую не используют.
\item Удаление недостижимого кода. Удаляет блоки без путей к ним. С помощью пропагации можно избавиться от ветвления с константным условием.
\item Встраивание функций. Заменяет вызов функции её телом. Это экономит время на вызов, но увеличивает размер. Компиляторы ограничиваются встраиванием малых функций.
\item Развертка циклов. Копирует тело цикла несколько раз, увеличивая тем самым объём кода, но сокращая число проверок. Данная техника способствует конвейеризации в силу меньшего ветвления.
\item Перемещение инвариантов. Выносит за пределы цикла выражения, не зависящие от переменной цикла и значение которого не меняется в течение выполнения всего цикла.
\item Слияние циклов. Комбинирует циклы с одинаковым количеством итераций и не зависящие друг от друга в один общий цикл. Это позволяет эффективнее использовать кэш.
\item Векторизация. Преобразует скалярные операции в векторные. Для массивов суммирование использует SIMD-инструкции.
\item Планирование инструкций. Переставляет команды для минимизации задержек. Зависимые инструкции разделяются независимыми.
\item Аллокация регистров. Распределяет переменные по регистрам. Алгоритм графового раскрашивания моделирует конфликты.
\item Глобальная нумерация. Значений выявляет эквивалентные выражения. SSA-форма их упрощает.
\item Хвостовая оптимизация рекурсии. Превращает рекурсию в цикл. В функциональных языках это предотвращает переполнение стека.
\item Оптимизации для кэша. Улучшают локальность. Перестановка циклов меняет порядок доступа к массиву.
\item Профильно-ориентированные техники. Используют данные выполнения. Компилятор переставляет ветви по вероятности.
\end{enumerate}

\subsection{Оптимизации компиляции функциональных языков и Lisp}

Основные методы оптимизаций функциональных языков программирования:

\begin{enumerate}
\item Оптимизация хвостовой рекурсии позволяет выполнять хвостовые вызовы без расширения стека вызовов и тем самым предотвращает переполнение стека при глубокой рекурсии. При компиляции хвостовой рекурсии вместо создания нового кадра стека происходит замена текущего кадра кадром следующего вызова, что превращает рекурсию в итерацию и тем самым экономит память.
\item Встраивание функций (inlining) -- это замена вызова функции её телом. Имеет различное влияние на производительность. Если функция вызывается один раз, то встраивание функции является очевидным решением, поскольку оно устраняет дополнительный вызов и упрощает выполнение кода. Если же функция вызывается много раз, то компилятор должен оценить, стоит ли делать подстановку функции. Это зависит от того, насколько эта подстановка повлияет на размер кода и скорость выполнения программы.
\item Structural Sharing -- это техника оптимизации неизменяемых структур данных, заключающаяся в том, что при создании новой версии структуры вместо полного копирования старой версии происходит её повторное использование. Например, после добавлении нового элемента в список этот элемент ссылается на старую часть списка. Это позволяет эффективнее использовать память, а также существенно улучшает скорость работы программ за счёт сокращения времени копирования структур.
\item Часто функциональные языки реализуют стратегию ленивых, или отложенных вычислений, согласно которой следует откладывать вычисления выражений до момента обращения к их значениям. Эта стратегия позволяет избежать ненужных вычислений, уменьшая нагрузку на процессор, и особенно полезна при работе с большими или бесконечными структурами данных. Однако ленивые вычисления не следует применять в языках, в которых присутствуют побочные эффекты, так как это может усложнить контроль над программой и привести к ошибкам. В чистых функциональных языках программирования отсутствуют побочные эффекты, поэтому для них использование отложенных вычислений является более естественным и предоставляет возможность писать чистый код.
\item Оптимизация копирования при записи (Copy-on-Write) является частным случаем ленивых вычислений и позволяет отложить копирование данных до момента их изменения. При дублировании структуры данных сначала передаётся ссылка на эти данные, вследствие чего данные становятся общими и разделяются между несколькими экземплярами. При попытке изменения одного из них происходит полное копирование исходных данных, после чего экземпляр получает собственную независимую копию данных. Данная техника оптимизирует память, производя копирование только тогда, когда оно становится необходимым.
\item Метод удаления мёртвого кода представляет собой удаление участков кода, которые никогда не выполнятся или не влияют на результат выполнения программы. Примером такого кода может служить функция, которая никогда не вызывается, ветка условия, которая никогда не вычисляется (при этом у компилятора есть достаточно информации, чтобы сделать такой вывод) или вычисление значения, которое нигде не используется. Данный метод уменьшает объём кода и время выполнения программы.
\item Техника удаления промежуточных структур снижает расход памяти путём объединения последовательных операций над структурами данных в общую операцию для избежания создания избыточных данных и их передачи между функциями. Эта оптимизация особенно актуальна в совокупности с ленивыми вычислениями, поскольку она компенсирует накладные расходы отложенных вычислений.
\item Удаление общих подвыражений заключается в обнаружении идентичных выражений для того, чтобы вычислить их один раз и затем повторно использовать полученный результат. Данный метод требует тщательный анализ программы на языке с побочными эффектами для точного выявления идентичных подвыражений. В силу отсутствия побочных эффектов в чистых функциональных языках программирования применение данной оптимизации особенно эффективно.
\item В качестве промежуточного представления кода при компиляции функциональных языков программирования зачастую используется стиль передачи продолжений (Continuation Passing Style), представляющий собой такой стиль написания кода, при котором каждая функция принимает дополнительный аргумент -- продолжение, то есть функцию, которая указывает, что выполнять дальше с результатом текущей функции. Вместо возврата значения напрямую функция вызывает это продолжение, передавая ему вычисленный результат. Преобразование в стиль передачи продолжений происходит рекурсивно по всей программе, вследствие чего управляющий поток программы становится явным, а контроль над порядком вычислений и возвращением значений сосредоточен в функциях-продолжениях.
\end{enumerate}

Особое место среди оптимизаций в Lisp занимает использование макросов, которые выступают важным инструментом для улучшения производительности и гибкости кода. Макросы разворачиваются уже на этапе компиляции, позволяя трансформировать и генерировать более эффективные программы. Это даёт возможность устранять избыточные вычисления и оптимизировать структуру кода ещё до его выполнения, обеспечивая тем самым уровень метапрограммирования, который является одной из ключевых особенностей Lisp и недоступен во многих других языках.

Несмотря на динамическую природу Lisp, современные компиляторы проводят довольно глубокий статический анализ. Например, в них активно используется добавление подсказок типов (type hints), что позволяет генерировать нативный код с меньшим количеством проверок во время выполнения. Такая стратегия уменьшает накладные расходы и ускоряет исполнение программ, сохраняя при этом гибкость динамической типизации.

Кроме того, важным аспектом оптимизации Lisp является эффективное управление памятью. Lisp-системы обычно применяют сборку мусора с поколенческим управлением, что существенно снижает задержки и повышает общую производительность системы. Компилятор, взаимодействующий со средой выполнения, может генерировать код, учитывающий особенности работы сборщика мусора, тем самым минимизируя накладные расходы, связанные с выделением и освобождением памяти.